\section{Discussion}
\label{sec:discussion}

\paragraph{Why both directions of dissociation matter.}
An agent that mentions an anchor in its summaries but does not
act on it occupies the most operationally dangerous quadrant: log
readouts look healthy while behaviour diverges. The inverse
quadrant (action lead) is less immediately dangerous but equally
invisible to verbal-only memory benchmarks, which would score the
cell as a recall failure and miss the fact that the directives
still weight the anchor. Both directions argue that a single-axis
verbal probe is insufficient to characterise the channel coupling
on tick-continuous agent loops. The pilot reports cells in both
directions on the same model: \textsf{fortified\_basin} runs
verbal-lead, \textsf{rugged\_highlands} runs action-lead, on
matched seeds and horizons.

\paragraph{Why scenario modulates direction.}
The 10-minute means separate by scenario rather than by seed,
which is the strongest cross-cell signal in the pilot. A mechanism
hypothesis is that the scenario's tile layout and initial
resource distribution determine which anchor key the LLM's prose
finds salient (the verbal half of the anchor) versus which the
LLM's policy continues to weight (the action half). In a basin
scenario the verbal channel surfaces the warehouse anchor sooner
than the policy re-aligns; in a highlands scenario the policy
weights a quarry/lumber anchor before the prose mentions it. We
frame this only as a hypothesis tied to two of nine short-horizon
cells crossing threshold; testing it directly requires a broader
scenario sweep.

\paragraph{Why horizon attenuates the gap.}
Magnitude shrinks at the longer horizon, with all three
60-minute scenario means inside a band of width $0.03$ around
zero. The substrate's action chain returns updated worker
positions and resource counts into each subsequent prompt, giving
the model many opportunities to observe its own dissociation and
correct. A cell whose 10-minute snapshot crosses threshold but
whose 60-minute snapshot does not is therefore not evidence of a
fading phenomenon: the dissociation appears at short horizons and
is most likely smoothed by the in-loop feedback by the time the
horizon closes. Without a finer time-resolution sweep we cannot
disentangle correction from natural drift, and we report the
attenuation as an observation.

\paragraph{Why the dual probe is not LongMemEval.}
A natural reviewer worry is that the dual probe's $\Delta$ might
reduce to a known long-context-recall failure. The external probe
rules this reading out for the same model at the same order of
context length: the pre-registered three-family external probe
(Section~\ref{sec:pilot:h2}) found every recovered cell scoring
perfectly, including the DeepSeek-v4-Flash cells whose substrate
$|\Delta|$ was largest on the original pilot. The substrate
dissociation is therefore distinguishable from the generic
retrieval-under-load failure that single-axis memory benchmarks
measure; a deployment monitor that only watched verbal-recall
health would have judged the dissociated cells fine.

\paragraph{Recommendation, with caveats.}
The dual probe is implementable on any structured-output agent
loop that emits both prose and a separately addressable action
surface. The substrate-side machinery here is one instantiation,
and its determinism makes the probe reproducible to a SHA-256
hash in the no-LLM control, but the probe itself does not depend
on the substrate. The recommendation we draw from the pilot is
that a paired action-grounded probe is worth co-monitoring with
verbal-only memory health checks on the kind of structured-output
agent loop we tested; whether the recommendation generalises to
production deployments outside the colony domain is the open
question Section~\ref{sec:limitations} flags as the gating
follow-on.
